\documentclass[11pt,letterpaper]{article}
% ==============================================================================
% Shared Preamble: Principles of Economics Modular Curriculum
% Core Colors: SlateBlue (#1E293B) & Terracotta (#C2410C)
% ==============================================================================

\usepackage[utf8]{inputenc}
\usepackage[T1]{fontenc}
\usepackage{lmodern}
\usepackage[margin=1in]{geometry}
\usepackage{amsmath,amssymb,amsfonts,amsthm}
\usepackage{booktabs}
\usepackage{microtype}
\usepackage{enumitem}
\usepackage{fancyhdr}
\usepackage{titlesec}
\usepackage{xcolor}
\usepackage{tikz}
\usepackage{pgfplots}
\usepackage{mdframed}
\usepackage[hidelinks,colorlinks=true,linkcolor=SlateBlue,citecolor=SlateBlue,urlcolor=Terracotta]{hyperref}

\pgfplotsset{compat=1.18}
\usetikzlibrary{arrows.meta,positioning,calc,patterns,decorations.pathreplacing}

% --- Color Palette Definitions ---
\definecolor{SlateBlue}{HTML}{1E293B}      % Primary Accent / Structure
\definecolor{Terracotta}{HTML}{C2410C}     % Secondary Accent / Emphasis
\definecolor{ForestGreen}{HTML}{15803D}    % Equilibrium / Positive Shocks
\definecolor{LightSlate}{HTML}{F8FAFC}     % Box Background Fill
\definecolor{BorderGray}{HTML}{CBD5E1}     % Neutral Borders
\definecolor{LightTerracotta}{HTML}{FFF7ED}% Cautionary / Principle Fill

% --- Section Heading Styling ---
\titleformat{\section}
  {\normalfont\Large\bfseries\color{SlateBlue}}{\thesection}{1em}{}[\titlerule]
\titleformat{\subsection}
  {\normalfont\large\bfseries\color{SlateBlue}}{\thesubsection}{1em}{}
\titleformat{\subsubsection}
  {\normalfont\normalsize\bfseries\color{Terracotta}}{\thesubsubsection}{1em}{}

% --- Header and Footer Configuration ---
\pagestyle{fancy}
\fancyhf{}
\renewcommand{\headrulewidth}{0.4pt}
\renewcommand{\footrulewidth}{0.4pt}
\fancyhead[L]{\small\scshape\color{SlateBlue}\leftmark}
\fancyhead[R]{\small\color{SlateBlue}Principles of Economics}
\fancyfoot[C]{\thepage}

% --- Custom Environments (mdframed) ---

% 1. Core Principle Box
\newmdenv[
  linecolor=Terracotta,
  linewidth=1.5pt,
  backgroundcolor=LightTerracotta,
  innertopmargin=10pt,
  innerbottommargin=10pt,
  innerrightmargin=12pt,
  innerleftmargin=12pt,
  roundcorner=4pt,
  skipabove=12pt,
  skipbelow=12pt
]{coreprinciple}

% 2. Empirical Case Study Box
\newmdenv[
  linecolor=SlateBlue,
  linewidth=1.5pt,
  backgroundcolor=LightSlate,
  innertopmargin=10pt,
  innerbottommargin=10pt,
  innerrightmargin=12pt,
  innerleftmargin=12pt,
  roundcorner=4pt,
  skipabove=12pt,
  skipbelow=12pt
]{casestudy}

% 3. Key Takeaways Box
\newmdenv[
  linecolor=BorderGray,
  linewidth=1pt,
  backgroundcolor=white,
  innertopmargin=10pt,
  innerbottommargin=10pt,
  innerrightmargin=12pt,
  innerleftmargin=12pt,
  roundcorner=4pt,
  skipabove=12pt,
  skipbelow=12pt
]{takeaways}

% Environment Title Helpers
\newcommand{\boxheader}[2]{{\noindent\bfseries\color{#1}\Large #2}\vspace{0.4em}\par}
\usepackage{tikz}
\usepackage{pgfplots}
\pgfplotsset{compat=1.18}
\usetikzlibrary{arrows.meta,positioning,calc,shapes.geometric,patterns,decorations.pathreplacing,intersections}

\tikzset{
  econ axis/.style={->, >=Stealth, thick, color=SlateBlue!90!black},
  econ grid/.style={thin, dashed, color=gray!30},
  econ guide/.style={dashed, color=gray!60, thin},
  demand curve/.style={very thick, color=SlateBlue},
  supply curve/.style={very thick, color=Terracotta},
  mb curve/.style={very thick, color=SlateBlue},
  mc curve/.style={very thick, color=Terracotta},
  econ point/.style={circle, fill=SlateBlue!90!black, inner sep=1.8pt},
  econ highlight/.style={fill=Terracotta!20, opacity=0.5},
  econ label/.style={font=\small\sffamily, align=center},
  econ annotation/.style={font=\footnotesize\itshape, color=SlateBlue!80!black}
}


\begin{document}

\title{\Huge\bfseries\color{SlateBlue} Chapter 11: Consumer Choice Theory: Utility and Indifference Curves}
\author{\large Principles of Microeconomics Curriculum}
\date{2026}
\maketitle

\thispagestyle{plain}

\begin{abstract}
This chapter provides a rigorous, pedagogical treatment of Chapter 11: Consumer Choice Theory: Utility and Indifference Curves, benchmarked to standard principles curriculum.
\end{abstract}

\tableofcontents
\vspace{1.5em}
\hrule
\vspace{1.5em}

\section{Learning Outcomes and Key Terms}
Students should represent a budget constraint, identify an affordable choice, and explain how prices and income shape a consumer's decision. These objectives support C2001-1 and C2001-3. Key terms include budget set, utility, indifference curve, marginal utility, marginal rate of substitution, substitution effect, and income effect.

\section{A Monthly Budget and Two Goods}
Suppose a student allocates a fixed monthly budget between meals ($X$) and transit trips ($Y$). The budget is limited, so buying more of one good means giving up some of the other. The consumer's choice depends on prices, income, and preferences; the model does not claim that every person calculates derivatives before buying lunch. It organizes the trade-offs implied by those constraints.

\section{The Budget Line: Geometric Description}
Place meals $X$ on the horizontal axis and transit trips $Y$ on the vertical axis. The budget line slopes downward. Its horizontal intercept is the number of meals affordable if all income is spent on meals; its vertical intercept is the number of transit trips affordable if all income is spent on transit. A rise in income shifts the budget line outward in parallel when prices are fixed. A change in one price pivots the line around the intercept for the other good.

The budget constraint is
\begin{equation}
P_xX+P_yY=I, \qquad Y=\frac{I}{P_y}-\frac{P_x}{P_y}X.
\end{equation}
The slope is $-P_x/P_y$: it is the number of transit trips that must be forgone to purchase one more meal while staying on the same budget line.

\section{Preferences, Utility, and Indifference Curves}
An indifference curve collects bundles that provide the same ordinal utility. A higher curve represents a preferred bundle when preferences are monotonic. With well-behaved preferences, indifference curves slope downward and are convex to the origin: as a person has more of $X$ relative to $Y$, they are generally willing to give up fewer units of $Y$ for another unit of $X$.

Marginal utility ($MU_x$) is the change in utility from a small increase in $X$, holding $Y$ fixed. The marginal rate of substitution is the absolute slope of an indifference curve:
\begin{equation}
MRS_{XY}=\frac{MU_x}{MU_y}.
\end{equation}
Utility numbers are a representation for ranking choices, not a literal measurement of happiness that can be compared across people.

\section{Choice at a Tangency: Worked Example}
Let income be $I=60$, with $P_x=4$ and $P_y=3$. The consumer has Cobb-Douglas utility $U(X,Y)=XY$. Since $MU_x=Y$ and $MU_y=X$, the tangency condition $MRS=P_x/P_y$ implies $Y/X=4/3$. Combining this with the budget constraint gives $X^*=7.5$ meals and $Y^*=10$ transit trips. The bundle exhausts the budget: $4(7.5)+3(10)=60$.

For an interior optimum with smooth, convex preferences:
\begin{equation}
\frac{MU_x}{MU_y}=\frac{P_x}{P_y} \quad\Longleftrightarrow\quad \frac{MU_x}{P_x}=\frac{MU_y}{P_y}.
\end{equation}
At a corner solution, the consumer chooses only one good; the tangency equality need not hold. Perfect substitutes, perfect complements, and nonstandard preferences can also change the shape of the solution.

\section{Price Changes: Substitution and Income Effects}
If the price of $X$ falls, the budget line pivots outward along the $X$ axis. The total change in the chosen amount of $X$ can be separated conceptually into a substitution effect and an income effect. The substitution effect reflects the change in relative prices while holding purchasing power conceptually constant; the income effect reflects the change in real purchasing power. For a normal good, both effects of a price decrease usually raise quantity demanded. For an inferior good, the income effect works in the opposite direction and may partly offset the substitution effect. A parallel budget shift is caused by a change in income at fixed prices; it is not the price-change income effect itself.

\begin{coreprinciple}
\boxheader{SlateBlue}{Clear It Up: The Price Change Pivots the Budget Line}
A change in $P_x$ changes the budget line's slope and rotates it around the $Y$ intercept. A change in income shifts the entire budget line in parallel when both prices remain fixed. Keep these two geometric changes distinct.
\end{coreprinciple}

\section{FRED Data and Real Purchasing Power}
FRED series CPIAUCSL is the Consumer Price Index for All Urban Consumers: All Items in U.S. City Average. It is monthly, seasonally adjusted, and indexed to 1982-1984=100; its September 2026 observation is 336.298. The CPI is a broad price index, not the exact cost of every household's basket. It can help convert a nominal dollar amount to a broad real-dollar comparison, but it does not identify an individual's preferences or their personal inflation rate.

\section{Check Your Understanding}
\begin{enumerate}
  \item With $I=60$, $P_x=4$, and $P_y=3$, what are the two intercepts of the budget line?
  \item Why does the budget line pivot rather than shift in parallel when only $P_x$ changes?
  \item Is utility measured in directly comparable units across people? Explain.
\end{enumerate}
Answers: (1) $X=15$ and $Y=20$. (2) The $X$ intercept changes while the $Y$ intercept remains $I/P_y$. (3) No; ordinal utility represents an individual's ranking of bundles.

\end{document}
